\section{Introduction}
\label{sec:intro}

Depth is what makes a neural network expressive
\citep{hornik1989multilayer,cybenko1989approximation,mhaskar2017deep,telgarsky2016benefits},
and it is also what makes it hard to train
\citep{glorot2010understanding,he2015delving,poole2016exponential,schoenholz2017deep}.
Two families of design decisions carry most of the weight in
keeping a deep stack trainable: the \emph{nonlinearity} applied between affine
maps, and the \emph{normalization} applied to the activations. The modern
defaults are well known --- ReLU~\citep{glorot2011deep,maas2013rectifier} or
GELU~\citep{hendrycks2016gaussian} for the nonlinearity, and
BatchNorm~\citep{ioffe2015batch} or LayerNorm~\citep{ba2016layer} for the
normalization --- but the space of alternatives is far from
exhausted, and the two decisions interact in ways that a single-axis ablation
cannot expose.

This paper takes as its starting point a specific ten-layer network in which
both decisions are made \emph{together}, by a single map:
\begin{equation}
\pi\bigl(f(z)\bigr) = \frac{f(z)}{\lVert f(z)\rVert_2},
\qquad
f(z) = z + \lambda z^{3},
\label{eq:projective-poly}
\end{equation}
where $z$ is the pre-activation of a layer and the $\ell_2$ norm is taken over
the feature dimension of each sample. We call this a \emph{projective
polynomial} activation: \emph{polynomial} because the nonlinearity is the
truncated cubic $z + \lambda z^3$ rather than a saturating or
rectifying function, and \emph{projective} because the output is pushed onto
the unit sphere, so each sample's feature vector is identified with a point of
projective space rather than of Euclidean space.

The construction is attractive on several grounds. It is
parameter-free (no learned scale or shift), it is batch-independent (unlike
BatchNorm, it does not couple samples within a mini-batch), it has no
train/eval discrepancy, and it provably bounds the activation magnitude,
which in a deep stack bounds the growth of the forward signal and, by the
chain rule, the growth of the backward signal. It also inherits the
projective-geometric intuition that motivates a body of work on the algebraic
geometry of linear networks, where the loss landscape of a deep stack is
studied through the projective geometry of its critical
points~\citep{xiuyi2021landscape,saxe2014exact,dauphin2014identifying}.

The difficulty with Eq.~\eqref{eq:projective-poly} as a scientific object is
that it bundles four separate claims:
\begin{enumerate}
\item[(C1)] a polynomial nonlinearity is a good nonlinearity;
\item[(C2)] an odd, unbounded-but-truncated polynomial is preferable to a
  rectifying or saturating one;
\item[(C3)] normalizing activations helps, and the specific normalization
  $\ell_2$ projection is preferable to centering-and-scaling alternatives;
\item[(C4)] \emph{depth} is where these differences matter, so the comparison
  must be made at depths where naive networks fail.
\end{enumerate}
Testing Eq.~\eqref{eq:projective-poly} as a monolith therefore tells us
whether \emph{this} design works, not \emph{which} of its ingredients is
responsible. The central methodological move of this paper is to pull the
construction apart: we implement the polynomial and the projection as
independent, independently switchable components, and then vary each against a
broad reference set.

\paragraph{What we compare.}
We organize the study around four reference families.
\begin{itemize}
\item \emph{Pointwise activations.} Identity (linear control), the classical
  saturating tanh, and the three non-saturating activations that current
  frontier language models actually deploy: ReLU~\citep{glorot2011deep}, GELU
  \citep{hendrycks2016gaussian} and SiLU~\citep{elfwing2018sigmoid}.
  We
  deliberately do not re-audit the historical zoo (sigmoid, LeakyReLU, ELU,
  SELU, Softplus, Mish, sinusoidal)~\citep{maas2013rectifier,clevert2015fast,klambauer2017self,misra2019mish,sitzmann2020implicit};
  none of those appears in a current
  large model, and our question is where a polynomial activation stands
  relative to what is deployed, not how the zoo ranks.
\item \emph{Gated linear units.} SwiGLU and GeGLU~\citep{shazeer2020glu,dauphin2017language} --- the shapes used in
  current large language models --- plus a polynomial-gated PolyGLU, which is
  the natural way to carry the paper's cubic into a gate. We include gating
  because ``the activation question'' in modern practice is often really a
  question about gating. Gated units contain two matrices per layer instead
  of one, so we match parameter count against the pointwise networks rather
  than comparing differently sized models.
\item \emph{Learnable polynomial activations.} A Legendre series, a Chebyshev
  series, and a KAN-style cubic B-spline~\citep{liu2024kan}. If a fixed cubic is useful, the
  natural question is whether letting the polynomial be learned improves on
  it; this family answers that question directly.
\item \emph{Normalization layers.} BatchNorm~\citep{ioffe2015batch},
  LayerNorm~\citep{ba2016layer}, RMSNorm~\citep{zhang2019root}, and the
  parameter-free $\ell_2$ projection, each placed both before and after the
  nonlinearity, and each crossed with both a rectified and a polynomial
  activation.
\end{itemize}

\paragraph{How we compare.}
All comparisons are made inside a single harness that isolates one factor at a
time, with the source network of Eq.~\eqref{eq:projective-poly} reproduced
bit-for-bit as the reference configuration (Sec.~\ref{sec:setup}). We vary
depth from 2 to 50 layers, polynomial order from 2 to 5, the polynomial
coefficient $\lambda$ over three decades, the $\ell_2$-projection placement,
the presence of a learnable activation, and the random seed. Alongside the
ten-class MNIST task we use eight synthetic datasets, including ones that are
deliberately hostile to a polynomial activation: a parity task, where the
label depends on a high-order interaction of binary features, and a
high-dimensional Gaussian task, where a low-dimensional signal is embedded in
a 50-dimensional ambient space.

\paragraph{Findings.}
Our main conclusions are the following.
\begin{itemize}
\item The projective polynomial network is \emph{not} competitive at equal
  parameter count: at depth ten the bare cubic $z + \lambda z^{3}$ lands almost
  four accuracy points \emph{below} ReLU, GELU and SiLU, and the parameter-free
  $\ell_2$ projection recovers a little over half of that deficit --- enough to
  bring the polynomial within two and a half points of the modern activations,
  but not enough to match them. Measured against the reference network's own
  depth, the polynomial is a reasonable but strictly dominated nonlinearity, and
  the paper's honest claim is about \emph{what the projection does}, not about
  the polynomial beating ReLU.
\item The performance of Eq.~\eqref{eq:projective-poly} comes
  \emph{predominantly from the projection, not the polynomial}:
  holding the polynomial fixed and removing the $\ell_2$ projection costs
  accuracy comparable to removing the nonlinearity's advantage over ReLU
  entirely, whereas holding the projection fixed and varying the polynomial
  changes accuracy only modestly.
\item At depth, the $\ell_2$ projection \emph{substitutes for} BatchNorm and
  LayerNorm: it recovers most of their benefit while using zero parameters and
  no batch statistics.
\item Making the polynomial learnable (Legendre, Chebyshev, B-spline) yields
  no consistent gain over the fixed cubic at equal depth and width: the
  learned bases are more expensive and no more accurate.
\item Gated units (SwiGLU, GeGLU) require normalization to train at depth at
  all, and once parameters are matched and normalization is present they are
  in the same band as the best pointwise activations; the large margin
  sometimes reported for gated units is a parameter-count effect, not an
  activation effect.
\item At depth ten on MNIST, a gated unit with no normalization fails to train
  at all, sitting at chance accuracy. The same unit trains once normalization
  is restored. This is the clearest single piece of evidence in our study
  that, for these architectures, \emph{normalization is not an optional
  refinement but a precondition for depth}.
\end{itemize}

The remainder of the paper is organized as follows.
Sec.~\ref{sec:related} positions the work. Sec.~\ref{sec:method} defines the
model and the activation and normalization families precisely.
Sec.~\ref{sec:setup} describes the experimental protocol.
Sec.~\ref{sec:results} reports results, and Sec.~\ref{sec:discussion}
discusses what they do and do not license. Sec.~\ref{sec:conclusion}
concludes.
